\documentclass[11pt,a4paper]{article}

% ===== PACKAGES =====
\usepackage[utf8]{inputenc}
\usepackage[T1]{fontenc}
\usepackage{geometry}
\usepackage{amsmath}
\usepackage{amssymb}
\usepackage{amsfonts}
\usepackage{graphicx}
\usepackage{hyperref}
\usepackage{fancyhdr}
\usepackage{titlesec}
\usepackage{enumitem}
\usepackage{float}
\usepackage{caption}
\usepackage{subcaption}
\usepackage{booktabs}
\usepackage{multirow}
\usepackage{xcolor}
\usepackage{tocloft}
\usepackage{listings}
\usepackage{algorithm}
\usepackage{algorithmic}
\usepackage{multicol}
\usepackage{array}

% Code listing style
\lstset{
    basicstyle=\ttfamily\small,
    breaklines=true,
    frame=single,
    backgroundcolor=\color{gray!10},
    keywordstyle=\color{blue},
    commentstyle=\color{green!60!black},
    stringstyle=\color{red!70!black}
}

% TOC styling
\renewcommand{\contentsname}{\Large\bfseries Contents}
\renewcommand{\cftsecfont}{\bfseries}
\renewcommand{\cftsecpagefont}{\bfseries}

\setlength{\cftbeforesecskip}{4pt}
\setlength{\cftbeforesubsecskip}{2pt}

\renewcommand{\cftdotsep}{1}

% ===== PAGE GEOMETRY =====
\geometry{
    a4paper,
    left=0.8in,
    right=0.8in,
    top=1in,
    bottom=1in
}

% ===== HYPERLINK SETUP =====
\hypersetup{
    colorlinks=true,
    linkcolor=blue,
    filecolor=magenta,      
    urlcolor=cyan,
    citecolor=blue,
    pdftitle={Velora Mobility Optimizer},
    pdfauthor={Hostel ID: 78},
}

% ===== SECTION FORMATTING =====
\titleformat{\section}
  {\normalfont\Large\bfseries}{\thesection}{1em}{}
\titleformat{\subsection}
  {\normalfont\large\bfseries}{\thesubsection}{1em}{}
\titleformat{\subsubsection}
  {\normalfont\normalsize\bfseries}{\thesubsubsection}{1em}{}

% ===== HEADER AND FOOTER =====
\pagestyle{fancy}
\fancyhf{}
\fancyhead[L]{Velora Mobility Optimizer}
\fancyhead[R]{Kriti '26}
\fancyfoot[C]{\thepage}
\renewcommand{\headrulewidth}{0.5pt}

% ===== TITLE PAGE INFO =====
\title{Velora Mobility Optimizer}
\author{Hostel ID: 78}
\date{\today}

% ===== DOCUMENT START =====
\begin{document}

% ===== CUSTOM TITLE PAGE =====
\begin{titlepage}
\thispagestyle{empty}
\centering

\vspace*{1cm}

{\large 3rd February 2026\par}

\vspace{2cm}

% LOGO
\includegraphics[width=0.40\textwidth]{image1.png}\par
\vspace{1.5cm}

{\Huge\bfseries Technical Report\par}
\vspace{0.5cm}

{\Large Kriti '26 --- Optimization Problem Statement\par}

\vspace{1.5cm}

\hrule height 1pt
\vspace{0.8cm}

{\Huge\bfseries Velora Mobility Optimizer\par}
\vspace{0.3cm}
{\Large\itshape Heterogeneous Vehicle Routing with Soft Time Windows\par}

\vspace{0.8cm}
\hrule height 1pt

\vspace{2.5cm}

{\Large Hostel ID:\par}
{\LARGE\bfseries 78\par}

\vfill

{\thepage}

\end{titlepage}

% ===== TABLE OF CONTENTS =====
\newpage
\tableofcontents
\thispagestyle{empty}
\newpage

% ===== ABSTRACT =====
\begin{abstract}
This report presents a comprehensive mathematical and algorithmic solution for the Velora corporate mobility optimization problem. The problem involves assigning employee transportation requests to a heterogeneous fleet of vehicles while minimizing operational cost and ensuring timely service, modeled as a \textbf{Heterogeneous Vehicle Routing Problem with Soft Time Windows (HVRPSTW)}.

Our solution implements a robust two-phase hybrid metaheuristic: (1) a \textbf{Multi-Start Greedy Constructive Heuristic} based on Solomon's I1 insertion criterion, and (2) a \textbf{Simulated Annealing (SA)} refinement phase employing Metropolis-Hastings acceptance criteria with multiple neighborhood operators (Relocate, Exchange, 2-Opt). The algorithm optimizes a weighted objective function combining monetary cost ($w_c=0.7$) and time penalty ($w_t=0.3$).

Experimental results on five distinct test cases demonstrate significant cost efficiencies, with savings ranging from \textbf{53.8\% to 77.0\%} compared to individual transportation baselines. The system adeptly handles complex constraints including vehicle capacity, priority-based time windows, and multi-vehicle type preferences, executing large-scale instances (100+ requests) in sub-second timeframes ($<$1s).
\end{abstract}

% ===== SECTION 1: INTRODUCTION =====
\section{Introduction}

\subsection{Problem Context}
Corporate employee transportation represents a significant operational cost and logistical challenge. The problem is characterized by static employee locations but dynamic daily demand, varying shift timings, and heterogeneous transport resources. The core challenge is to partition the set of employee requests into vehicle routes such that all constraints are satisfied and the total operational cost is minimized.

\subsection{Objective}
Velora's objective is to minimize a composite cost function:
\begin{itemize}
    \item \textbf{Direct Operating Costs}: Fuel and vehicle usage costs (weighted at 0.7).
    \item \textbf{Service Quality Penalties}: Late arrivals, early waits, and preference violations (weighted at 0.3).
\end{itemize}

This must be achieved while strictly adhering to:
\begin{enumerate}
    \item \textbf{Capacity Constraints}: Vehicle passenger limits.
    \item \textbf{Precedence Constraints}: Pickup must occur before dropoff.
    \item \textbf{Hard Time Windows}: Maximum allowable delay thresholds.
\end{enumerate}

\subsection{Problem Classification}
This problem is formally classified as \textbf{HVRPSTW} (Heterogeneous Vehicle Routing Problem with Soft Time Windows). It generalizes the classical VRP by introducing:
\begin{itemize}
    \item \textbf{Heterogeneity}: $K$ vehicle types, each with distinct capacity $Q_k$ and cost $C_k$.
    \item \textbf{Soft Time Windows}: Violation of optimal time windows $[E_i, L_i]$ is penalized rather than strictly forbidden (up to a tolerance $\Delta$).
    \item \textbf{Many-to-One Flow}: Multiple pickups converging to common drop-offs (or distinct drop-offs in the general case).
\end{itemize}

% ===== SECTION 2: PROBLEM FORMULATION =====
\section{Problem Formulation}

\subsection{Mathematical Model}
We define the problem on a directed graph $G = (V, A)$, where $V = \{0\} \cup P \cup D$ is the set of vertices (0 is depot, $P$ pickups, $D$ drop-offs), and $A$ is the set of arcs.

\subsubsection{Sets and Parameters}
\begin{itemize}
    \item $N = \{1, \ldots, n\}$: Set of requests. Each request $i \in N$ has a pickup node $p_i \in P$ and dropoff node $d_i \in D$.
    \item $K$: Set of available vehicles.
    \item $d_{ij}$: Euclidean distance between nodes $i$ and $j$.
    \item $c_{ij}^k$: Cost of traversing arc $(i,j)$ with vehicle $k$.
    \item $t_{ij}^k$: Travel time on arc $(i,j)$ with vehicle $k$.
    \item $Q_k$: Capacity of vehicle $k$.
    \item $[E_i, L_i]$: Target time window for request $i$.
\end{itemize}

\subsubsection{Decision Variables}
\begin{itemize}
    \item $x_{ijk} \in \{0, 1\}$: Binary variable, equal to 1 if vehicle $k$ traverses arc $(i, j)$.
    \item $y_{ik} \in \{0, 1\}$: Binary variable, equal to 1 if request $i$ is served by vehicle $k$.
    \item $T_{ik}$: Service start time at node $i$ by vehicle $k$.
    \item $L_{ik}$: Accumulation load variable after visiting node $i$.
\end{itemize}

\subsubsection{Objective Function}
The objective is to minimize the total generalized cost $Z$:

\begin{equation}
\min Z = \alpha \sum_{k \in K}\sum_{(i,j) \in A} c_{ij}^k x_{ijk} + \beta \sum_{i \in N} \text{Penalty}(T_{drop_i})
\end{equation}

Where $\alpha=0.7, \beta=0.3$. The penalty function is piecewise linear:
\begin{equation}
\text{Penalty}(t) = 
\begin{cases} 
0 & \text{if } t \le L_i \\
w_i \cdot (t - L_i) & \text{if } L_i < t \le L_i + \Delta_{max} \\
\infty & \text{if } t > L_i + \Delta_{max}
\end{cases}
\end{equation}
Where $w_i$ is the priority weight for employee $i$.

\subsubsection{Constraints}
\begin{align}
\sum_{k \in K} \sum_{j \in V} x_{ijk} &= 1 & \forall i \in P \quad \text{(Each request picked up once)} \\
\sum_{j \in V} x_{p_i, j, k} - \sum_{j \in V} x_{j, d_i, k} &= 0 & \forall i \in N, k \in K \quad \text{(Flow conservation)} \\
T_{p_i, k} + t_{p_i, d_i}^k &\le T_{d_i, k} & \forall i \in N, k \in K \quad \text{(Precedence)} \\
\sum_{i \in N} q_i y_{ik} &\le Q_k & \forall k \in K \quad \text{(Capacity)}
\end{align}

% ===== SECTION 3: ALGORITHM DESIGN =====
\section{Algorithm Design}

To address the NP-hardness of HVRPSTW, we implement a two-stage metaheuristic.

\subsection{Phase 1: Multi-Start Constructive Heuristic}
This phase generates a set of feasible initial solutions using a randomized greedy insertion strategy (Solomon's I1).

\subsubsection{Initialization}
Requests are sorted by priority ($P_1 > P_5$) and then by time urgency (earliest deadline first). We apply random perturbations to this order in each of the 10 restart iterations to explore different parts of the solution space.

\subsubsection{Insertion Criteria}
For a request $u$ and a vehicle $k$ with route $R_k = (v_0, v_1, \ldots, v_m)$, we evaluate all possible insertion pairs $(i, j)$ where $0 \le i < j \le m+1$. The best insertion position minimizes the weighted insertion cost:

\begin{equation}
\Delta C(u, k) = \min_{i,j} \left[ \alpha \cdot \Delta \text{Dist}_{i,j} + \beta \cdot \Delta \text{Time}_{i,j} + \Delta \text{Penalty} \right]
\end{equation}

Where $\Delta \text{Dist}_{i,j}$ represents the detours added by inserting pickup at index $i$ and dropoff at index $j$.

\subsection{Phase 2: Simulated Annealing (SA)}
The best solution from Phase 1 is refined using Simulated Annealing, a probabilistic technique for approximating the global optimum of a given function.

\subsubsection{Theoretical Basis}
SA mimics the physical cooling of solids. The state space is traversed using a Markov Chain where a transition from state $S$ to $S'$ with energy (cost) difference $\Delta E = E(S') - E(S)$ is accepted with probability:

\begin{equation}
P(\text{accept}) = \begin{cases} 
1 & \text{if } \Delta E < 0 \\
e^{-\Delta E / T} & \text{if } \Delta E \ge 0 
\end{cases}
\end{equation}

This property allows the algorithm to escape local minima (hill-climbing) early in the search when temperature $T$ is high.

\subsubsection{Cooling Schedule}
\begin{itemize}
    \item \textbf{Initial Temperature ($T_0$)}: 1000.0
    \item \textbf{Cooling Rate ($\alpha$)}: 0.995 (Geometric cooling scheme $T_{k+1} = \alpha T_k$)
    \item \textbf{Stopping Criteria}: $T < 0.1$ or Max Iterations without improvement $> 2000$.
\end{itemize}

\subsubsection{Neighborhood Operators}
We implement three mutation operators to generate neighbor $S'$:
\begin{enumerate}
    \item \textbf{Relocate ($O(1)$)}: Moves a single request $r$ from route $R_a$ to route $R_b$.
    \item \textbf{Exchange ($O(1)$)}: Swaps request $r_1$ in $R_a$ with request $r_2$ in $R_b$.
    \item \textbf{2-Opt ($O(n^2)$)}: Time-reversible heuristic. Reverses the sequence of stops between indices $i$ and $j$ in a single route. This untangles crossing paths.
\end{enumerate}

% ===== SECTION 4: SYSTEM ARCHITECTURE =====
\section{System Architecture}

\subsection{Software Stack}
\begin{center}
\begin{tabular}{|l|l|l|}
\hline
\textbf{Component} & \textbf{Language} & \textbf{Role} \\
\hline
Solver Core & C++17 & High-performance combinatorial optimization \\
Data Parser & Python 3.9 & Excel ingestion, data validation \\
Reporting & Python 3.9 & Report generation, cost analysis \\
Build System & CMake & Cross-platform compilation \\
\hline
\end{tabular}
\end{center}

\subsection{Module Interaction}
Executed via `run_pipeline.sh`:
\begin{enumerate}
    \item `excel_to_json.py`: Parses raw requests, geocodes addresses if API available, else validates lat/long.
    \item `velora_solver`: Reads JSON, executes SA algorithm, writes optimized routes to JSON.
    \item `json_to_report.py`: Calculates financial metrics (CTC) and generates the readable report.
\end{enumerate}

\subsection{Distance Calculation Logic}
All modules share a unified `MapDistance` class to ensure consistency:
\begin{itemize}
    \item \textbf{Primary}: Google Maps Distance Matrix API (Travel time accuracy).
    \item \textbf{Fallback}: Adjusted Haversine Formula:
    \begin{equation}
    D_{road} \approx D_{haversine} \times 1.4
    \end{equation}
    The 1.4 factor acts as a statistically derived "tortuosity factor" to approximate urban road networks.
\end{itemize}

% ===== SECTION 5: EXPERIMENTAL RESULTS =====
\section{Experimental Results}

We evaluated the system on five test cases (TC01-TC05) representing increasing complexity.

\subsection{Summary Metrics}
\begin{center}
\begin{tabular}{lcccccc}
\toprule
\textbf{Case} & \textbf{Reqs} & \textbf{Base Cost (₹)} & \textbf{Opt Cost (₹)} & \textbf{Savings} & \textbf{Vehicles} & \textbf{Runtime} \\
\midrule
TC01 & 8 & 3,200 & 920 & \textbf{71.2\%} & 2 & 85ms \\
TC02 & 12 & 4,945 & 1,273 & \textbf{74.3\%} & 3 & 120ms \\
TC03 & 15 & 6,450 & 1,482 & \textbf{77.0\%} & 3 & 150ms \\
TC04 & 10 & 4,420 & 2,044 & \textbf{53.8\%} & 4 & 180ms \\
TC05 & 100 & 45,000 & 18,320 & \textbf{59.3\%} & 22 & 890ms \\
\bottomrule
\end{tabular}
\end{center}

\subsection{Detailed Analysis: TC01 (Basic)}
\begin{itemize}
    \item \textbf{Scenario}: 8 employees, 3 allocated vehicles.
    \item \textbf{Outcome}: Only 2 vehicles needed (V01, V02). V03 remained unused.
    \item \textbf{Efficiency}: Load factors of 100\% (3/3) and 100\% (4/4) achieved on active routes.
\end{itemize}

\subsection{Detailed Analysis: TC04 (Stress Test)}
TC04 introduced an "impossible" request (E10) where travel time exceeded the time window.
\begin{itemize}
    \item \textbf{Constraint Handling}: The solver correctly identified the impossibility.
    \item \textbf{Fallback}: Instead of declaring infeasibility, it applied the "Forced Assignment" strategy, serving E10 but accumulating a penalty cost.
    \item \textbf{Cost}: Higher optimization cost due to the penalties, but still 53.8\% cheaper than baseline.
\end{itemize}

% ===== SECTION 6: CHALLENGES FACED =====
\section{Challenges Faced}

\subsection{1. Pickup-Dropoff Precedence in Local Search}
\textbf{Challenge}: The standard 2-opt operator blindly reverses route segments. If a pickup and its corresponding drop-off are both within the reversed segment, their order is swapped (Dropoff $\rightarrow$ Pickup), violating physical reality.
\textbf{Solution}: Modified `isRouteFeasible()` function to perform an $O(n)$ scan of every route after any permutation, verifying that for all requests $i$, index($pickup_i$) $<$ index($dropoff_i$).

\subsection{2. Inconsistent Distance Metrics}
\textbf{Challenge}: Initially, the solver used Haversine distance while the report generator used Euclidean distance, leading to cost discrepancies ($D \neq D'$).
\textbf{Solution}: Centralized distance logic into a shared C++ header and Python module that both implement the exact same math: Haversine $\times$ 1.4 Tortuosity Factor.

\subsection{3. Convergence of Simulated Annealing}
\textbf{Challenge}: With high penalties for time windows, the "energy landscape" became very rugged. The algorithm would get stuck in local minima where moving any request incurred massive penalties.
\textbf{Solution}: Tuned the cooling schedule. Increased initial temperature $T_0$ to 1000 to allow more "uphill" moves early on, enabling the solver to jump out of deep local minima.

% ===== SECTION 7: CONCLUSIONS =====
\section{Conclusion}

The Velora Mobility Optimizer successfully demonstrates that heuristic optimization can deliver massive efficiencies in corporate transportation. By moving from static/individual routing to dynamic shared routing:
\begin{enumerate}
    \item \textbf{Cost Reduction}: Consistent savings of $>$50\% are achievable.
    \item \textbf{Fleet Utilization}: Reduction in required vehicles by 20-30\%.
    \item \textbf{Scalability}: The $O(N^2)$ complexity of the core heuristic ensures the system remains responsive even as employee count grows to hundreds.
\end{enumerate}

The implementation of Simulated Annealing ensures that given sufficient compute time, the solution quality approaches the global optimum, while the Constructive Heuristic guarantees a "good enough" solution is available instantly.

% ===== SECTION 8: RESOURCES =====
\section{Resources and Code Availability}

\subsection{Repository}
The complete source code, test data, and build scripts are hosted on GitHub:
\begin{quote}
\url{https://github.com/badri41/velora-mobility-optimizer}
\end{quote}

\subsection{Project Structure}
\begin{description}
    \item[solver/] Core C++17 optimization engine.
    \item[parser/] Python modules for Excel/JSON transformation.
    \item[scripts/] Shell scripts for pipeline orchestration.
    \item[data/] Contains all 5 Test Cases (Input Excel/JSON and Output Reports).
\end{description}

% ===== SECTION 9: MODULES AND ANSWERS =====
\section{Modules and Answers}

\begin{multicols}{2}
\noindent
\textbf{Module 1}\\
Question 1 \dotfill 4\\
Question 2 \dotfill 14\\
Question 3 \dotfill 11\\

\noindent
\textbf{Module 2}\\
Question 1 \dotfill 9\\
Question 2 \dotfill 10\\

\noindent
\textbf{Module 3}\\
Question 1 \dotfill 5\\
Question 2 \dotfill 7\\

\noindent
\textbf{Module 4}\\
Question 1 \dotfill 28\\

\noindent
\textbf{Module 5}\\
Question 1 \dotfill 14\\

\noindent
\textbf{Module 6}\\
Question 1 \dotfill 15\\
Question 2 \dotfill 20\\
Question 3 \dotfill 19\\

\noindent
\textbf{Module 7}\\
Question 1 \dotfill 26\\
Question 2 \dotfill 27\\

\noindent
\textbf{Module 8}\\
Question 1 \dotfill 22\\
Question 2 \dotfill 24\\

\noindent
\textbf{Module 9}\\
Question 1 \dotfill 29\\
Question 2 \dotfill 33\\
\end{multicols}

% ===== REFERENCES =====
\begin{thebibliography}{99}

\bibitem{solomon1987}
Solomon, M.M.,
\textit{Algorithms for the Vehicle Routing and Scheduling Problems with Time Window Constraints},
Operations Research, vol. 35, no. 2, pp. 254-265, 1987.

\bibitem{kirkpatrick1983}
Kirkpatrick, S., Gelatt, C.D., and Vecchi, M.P.,
\textit{Optimization by Simulated Annealing},
Science, vol. 220, no. 4598, pp. 671-680, 1983.

\bibitem{toth2014}
Toth, P., and Vigo, D.,
\textit{Vehicle Routing: Problems, Methods, and Applications},
SIAM, 2014.

\bibitem{laporte2009}
Laporte, G.,
\textit{Fifty Years of Vehicle Routing},
Transportation Science, vol. 43, no. 4, pp. 408-416, 2009.

\end{thebibliography}

\end{document}
